% Gabarit pandoc des corrigés publiés (brut et guidé). Contenu sans marque ni nom de site.
% Compilation : lualatex (latexmk -lualatex).
% Le Markdown reste la source unique : ne pas éditer ce fichier .tex à la main.
\documentclass[11pt,a4paper]{article}

% --- Moteur, polices (Latin Modern) et langue ---------------------------------
\usepackage{iftex}
\ifLuaTeX\else\errmessage{Ce document se compile avec lualatex}\fi
\usepackage{amsmath,amssymb,mathtools}
\usepackage{fontspec}
% mathrm=sym : \mathrm prend ses glyphes dans la police mathématique (ℓ, μ, Ω), comme KaTeX.
\usepackage[math-style=TeX,bold-style=TeX,mathrm=sym]{unicode-math}
% Repli pour les glyphes absents de Latin Modern (indices et exposants Unicode,
% flèches, filets des schémas en texte, symboles) : DejaVu, puis Latin Modern Math
% (ℓ, lettres mathématiques Unicode des légendes).
\directlua{
  luaotfload.add_fallback("repliTexte", {"DejaVuSerif:mode=harf;", "DejaVuSans:mode=harf;", "file:latinmodern-math.otf:mode=harf;"})
  luaotfload.add_fallback("repliMono", {"DejaVuSansMono:mode=harf;", "DejaVuSans:mode=harf;", "file:latinmodern-math.otf:mode=harf;"})
}
\setmainfont{Latin Modern Roman}[RawFeature={fallback=repliTexte}]
\setsansfont{Latin Modern Sans}[RawFeature={fallback=repliTexte}]
\setmonofont{Latin Modern Mono}[RawFeature={fallback=repliMono}]
\setmathfont{Latin Modern Math}
\usepackage[french]{babel}
\frenchsetup{SmallCapsFigTabCaptions=false}

% --- Mise en page --------------------------------------------------------------
\usepackage[a4paper,margin=2.2cm,headheight=14pt,footskip=1.2cm]{geometry}
\usepackage{microtype}
\setlength{\emergencystretch}{4em}
\tolerance=1500
\setlength{\parindent}{0pt}
\setlength{\parskip}{0.55\baselineskip plus 2pt minus 1pt}
\usepackage{xcolor}
% Gris nommés : babel-french rend « ! » actif dans le document.
\definecolor{gris2}{gray}{0.98}
\definecolor{gris3}{gray}{0.97}
\definecolor{gris5}{gray}{0.95}
\definecolor{gris35}{gray}{0.65}
\definecolor{gris40}{gray}{0.60}
\definecolor{gris45}{gray}{0.55}
\definecolor{gris50}{gray}{0.50}
\definecolor{gris55}{gray}{0.45}
\definecolor{gris60}{gray}{0.40}
\definecolor{gris65}{gray}{0.35}
\definecolor{gris70}{gray}{0.30}
\definecolor{gris80}{gray}{0.20}
\definecolor{gris85}{gray}{0.15}
\usepackage{graphicx}
\usepackage{float}
\usepackage{caption}
\captionsetup{font=small,labelformat=empty,justification=centering}
\usepackage{booktabs,longtable,array,multirow,calc}
\usepackage{siunitx}
\sisetup{locale=FR}
\usepackage[version=4]{mhchem}
\usepackage{listings}
\usepackage{fancyhdr}
\usepackage{lastpage}
\usepackage[most]{tcolorbox}
\usepackage{enumitem}
\setlist{nosep,leftmargin=*,topsep=2pt}
\usepackage[hidelinks,unicode]{hyperref}
\usepackage{bookmark}
\hypersetup{pdftitle={Mines-Ponts MP-MPI-PC-PSI — Langue vivante 2026},pdfauthor={Rédigé par IA},pdfcreator={LaTeX},pdfsubject={Corrigé guidé}}
\usepackage{xurl}

% --- Pandoc : listes, tableaux, figures ----------------------------------------
\providecommand{\passthrough}[1]{#1}
\providecommand{\tightlist}{\setlength{\itemsep}{0pt}\setlength{\parskip}{0pt}}

\makeatletter
\def\maxwidth{\ifdim\Gin@nat@width>\linewidth\linewidth\else\Gin@nat@width\fi}
\def\maxheight{\ifdim\Gin@nat@height>0.55\textheight 0.55\textheight\else\Gin@nat@height\fi}
\makeatother
\setkeys{Gin}{width=\maxwidth,height=\maxheight,keepaspectratio}
\setcounter{secnumdepth}{-\maxdimen}
\providecommand{\pandocbounded}[1]{#1}
\usepackage{longtable,booktabs,array}

% --- Code (listings) -----------------------------------------------------------
\lstset{
  basicstyle=\ttfamily\small,
  columns=fullflexible,
  keepspaces=true,
  breaklines=true,
  breakatwhitespace=false,
  postbreak=\mbox{\textcolor{gris50}{$\hookrightarrow$}\space},
  frame=single,
  rulecolor=\color{gris35},
  framesep=4pt,
  xleftmargin=5pt,xrightmargin=5pt,
  backgroundcolor=\color{gris3},
  keywordstyle=\bfseries,
  commentstyle=\color{gris65},
  stringstyle=\color{gris80},
  showstringspaces=false,
  tabsize=4,
  extendedchars=true,
  aboveskip=6pt,belowskip=6pt,
}

% --- Macros des corrigés (KaTeX ↔ LaTeX) ---------------------------------------
% Les corrigés sont écrits pour KaTeX : on définit ici ce que KaTeX connaît et
% que LaTeX n'a pas, sans jamais retoucher le Markdown.
\providecommand{\R}{\mathbb{R}}
\providecommand{\N}{\mathbb{N}}
\providecommand{\Z}{\mathbb{Z}}
\providecommand{\Q}{\mathbb{Q}}
\providecommand{\C}{\mathbb{C}}
\providecommand{\dd}{\mathrm{d}}
\providecommand{\llbracket}{\lBrack}
\providecommand{\rrbracket}{\rBrack}
\providecommand{\lBrace}{\lbrace}
\providecommand{\rBrace}{\rbrace}
\providecommand{\Box}{\square}
\providecommand{\Diamond}{\lozenge}
\providecommand{\varnothing}{\emptyset}
\providecommand{\argmax}{\operatorname*{arg\,max}}
\providecommand{\argmin}{\operatorname*{arg\,min}}
\renewcommand{\leq}{\leqslant}
\renewcommand{\geq}{\geqslant}

% Équation hors texte : réduite à la largeur de la ligne si elle déborde.
\newsavebox{\boiteequation}
\newcommand{\equationajustee}[1]{%
  \sbox{\boiteequation}{$\displaystyle #1$}%
  \ifdim\wd\boiteequation>\linewidth
    \[\resizebox{\linewidth}{!}{\usebox{\boiteequation}}\]%
  \else
    \[#1\]%
  \fi}

% --- Encadrés pédagogiques (sobres, lisibles en noir et blanc) -----------------
\tcbset{
  encadre/.style={
    enhanced,breakable,sharp corners,
    colback=white,colframe=gris55,
    boxrule=0.4pt,leftrule=2.2pt,
    left=6pt,right=6pt,top=3pt,bottom=3pt,
    before skip=6pt,after skip=6pt,
    fonttitle=\bfseries\small,coltitle=black,
    attach title to upper,after title={\quad},
    parbox=false,
  }
}
\newtcolorbox{indication}[1]{encadre,colframe=gris45,colback=gris2,title={Indication~#1}}
\newtcolorbox{methode}{encadre,colframe=gris70,title={Méthode}}
\newtcolorbox{piege}{encadre,colframe=gris85,boxrule=0.8pt,leftrule=2.2pt,title={Piège}}
\newtcolorbox{jury}{encadre,colframe=gris60,colback=gris5,title={Jury}}
\newtcolorbox{encadrenote}{encadre,colframe=gris40}

% --- Titres de section ---------------------------------------------------------
\usepackage{titlesec}
\titleformat{\section}{\Large\bfseries}{}{0pt}{}
\titleformat{\subsection}{\large\bfseries}{}{0pt}{}
\titleformat{\subsubsection}{\normalsize\bfseries}{}{0pt}{}
\titlespacing*{\section}{0pt}{14pt plus 4pt minus 2pt}{6pt}
\titlespacing*{\subsection}{0pt}{10pt plus 3pt minus 2pt}{4pt}
\titlespacing*{\subsubsection}{0pt}{8pt plus 2pt minus 2pt}{2pt}

% --- En-tête et pied de page ---------------------------------------------------
\pagestyle{fancy}
\fancyhf{}
\fancyhead[L]{\small\textsc{Corrigé guidé} · Mines-Ponts MP-MPI-PC-PSI}
\fancyhead[R]{\small Langue vivante 2026}
\fancyfoot[L]{\footnotesize Corrigé rédigé par IA, non officiel}
\fancyfoot[R]{\footnotesize \thepage\,/\,\pageref*{LastPage}}
\renewcommand{\headrulewidth}{0.4pt}
\renewcommand{\footrulewidth}{0pt}
\fancypagestyle{plain}{\fancyhf{}\fancyfoot[R]{\footnotesize \thepage\,/\,\pageref*{LastPage}}\fancyfoot[L]{\footnotesize Corrigé rédigé par IA, non officiel}\renewcommand{\headrulewidth}{0pt}}

\begin{document}
\thispagestyle{plain}
\begin{center}
  {\small\scshape Corrigé guidé}\par\smallskip
  {\LARGE\bfseries Mines-Ponts MP-MPI-PC-PSI — Langue vivante 2026\par}\medskip
  {\large Mines-Ponts · MP-MPI-PC-PSI · Langue vivante · 2026\par}
\end{center}
\vspace{2pt}\hrule\vspace{8pt}

\emph{Corrigé rédigé par IA, non officiel. Signalez toute erreur.}

\hypertarget{pruxe9sentation-du-sujet}{%
\section{Présentation du sujet}\label{pruxe9sentation-du-sujet}}

\textbf{Format.} L'épreuve dure 1 h 30 et est commune aux filières MP, MPI, PC et PSI. Dictionnaires et appareils sont interdits. Chaque langue comporte trois exercices :

\begin{longtable}[]{@{}
  >{\raggedright\arraybackslash}p{(\columnwidth - 6\tabcolsep) * \real{0.1655}}
  >{\raggedright\arraybackslash}p{(\columnwidth - 6\tabcolsep) * \real{0.3846}}
  >{\raggedright\arraybackslash}p{(\columnwidth - 6\tabcolsep) * \real{0.2719}}
  >{\raggedright\arraybackslash}p{(\columnwidth - 6\tabcolsep) * \real{0.1480}}@{}}
\toprule\noalign{}
\begin{minipage}[b]{\linewidth}\raggedright
Exercice
\end{minipage} & \begin{minipage}[b]{\linewidth}\raggedright
Nature
\end{minipage} & \begin{minipage}[b]{\linewidth}\raggedright
Longueur
\end{minipage} & \begin{minipage}[b]{\linewidth}\raggedright
Barème
\end{minipage} \\
\midrule\noalign{}
\endhead
\bottomrule\noalign{}
\endlastfoot
Question 1 & compréhension de l'article, reformulée & 80 mots ± 10 \% (72 à 88) & 4 points \\
Question 2 & expression personnelle argumentée & 180 mots ± 10 \% (162 à 198) & 8 points \\
Thème & traduction d'un extrait de roman français contemporain & texte entier & 8 points \\
\end{longtable}

Le candidat indique le nombre de mots de chaque réponse et écrit en tête de copie la langue de composition, qui doit être celle de son dossier d'inscription. Il ne traduit ni les références ni les titres du thème, et il écrit une ligne sur deux. La version arabe est la seule à chiffrer les bornes (72-88 et 162-198) ; elles sont la traduction exacte du « ± 10 \% ». Le candidat ne compose que dans une langue, mais ce corrigé traite les six, soit 18 exercices numérotés de 1 à 18.

\textbf{Les six dossiers.}

\begin{longtable}[]{@{}
  >{\raggedright\arraybackslash}p{(\columnwidth - 6\tabcolsep) * \real{0.0710}}
  >{\raggedright\arraybackslash}p{(\columnwidth - 6\tabcolsep) * \real{0.3388}}
  >{\raggedright\arraybackslash}p{(\columnwidth - 6\tabcolsep) * \real{0.2874}}
  >{\raggedright\arraybackslash}p{(\columnwidth - 6\tabcolsep) * \real{0.2728}}@{}}
\toprule\noalign{}
\begin{minipage}[b]{\linewidth}\raggedright
Langue
\end{minipage} & \begin{minipage}[b]{\linewidth}\raggedright
Article (expression écrite)
\end{minipage} & \begin{minipage}[b]{\linewidth}\raggedright
Ce que demandent les questions
\end{minipage} & \begin{minipage}[b]{\linewidth}\raggedright
Thème
\end{minipage} \\
\midrule\noalign{}
\endhead
\bottomrule\noalign{}
\endlastfoot
Allemand & \emph{Tagesschau}, 2025 : les \emph{Serious Games} (« The Darkest Files », enquêtes sur des crimes nazis ; « Das Ilios-Experiment », tolérance), intérêt pédagogique et limites (\emph{kein Allheilmittel}) & Q1 : possibilités offertes par ces jeux ; Q2 : l'école doit-elle aussi transmettre des valeurs ? avec exemples & Le Clézio, \emph{Bitna} : une tante demande à la narratrice de sermonner une adolescente qui fume \\
Anglais & \emph{The Guardian}, 2025 : déclin d'Asos après le Covid (stock d'un milliard de livres, Shein, Vinted, coûts) & Q1 : pourquoi Asos peine à rester compétitif, « in your own words » ; Q2 : le prix doit-il être le seul critère d'achat ? & Foenkinos, \emph{Vers la beauté} : une adolescente dépressive, sa mère, un médecin \\
Arabe & \emph{Al-Araby Al-Jadeed}, 2025 : le mouvement marocain « Gen Z », fin des grandes idéologies & Q1 : causes des protestations selon le texte, avec ses propres mots ; Q2 : mutation durable ou phénomène passager ? exemples arabes & Adimi, \emph{La Joie ennemie} : des parents algériens bannissent l'arabe de la maison (texte non vocalisé) \\
Espagnol & \emph{El País}, 2025 : 0,8 \% des jeunes Gitans diplômés du supérieur & Q1 : obstacles identifiés ; Q2 : mesures d'intégration des minorités en Espagne et en Amérique latine, avec exemples hispaniques & Benzine, \emph{Les silences des pères} : un mineur marocain raconte la descente au fond (contextualisation à ne pas traduire) \\
Italien & \emph{Corriere della Sera}, 2025 : canicule précoce, \emph{bollino rosso}, Méditerranée surchauffée & Q1 : illustrer le problème ; Q2 : la crise climatique sera-t-elle résolue, et comment ? & Martin, \emph{Mes fous} : un narrateur dépendant de son portable \\
Russe & \emph{Школа Жизни}, 2025 : codes vestimentaires et argot de la culture jeune (hip-hop, « лол », « рофл ») & Q1 : style de la jeunesse ; Q2 : dans quelle mesure vous y reconnaissez-vous ? & Davrichewy, \emph{L'autre Joseph} : une grand-mère, un arrière-grand-père absent \\
\end{longtable}

\textbf{Thème de la session.} Les articles portent sur des faits de société récents et chiffrés. Les thèmes sont des extraits de romans français parus entre 2016 et 2025. Ils mêlent récit au passé ou au présent de narration et dialogue familier. Leur difficulté tient moins au vocabulaire qu'aux temps, à l'aspect, au mode (subjonctif, discours indirect) et au rendu idiomatique de tours oraux.

\textbf{Ce qu'il ne faut pas laisser passer.} Il ne faut sortir ni de la fourchette de mots ni oublier d'indiquer le nombre de mots : ce sont deux consignes explicites. En question 1, on répond directement, sans introduction, et on reformule : le recopiage se voit tout de suite. En question 2, on écrit une courte introduction qui situe le sujet dans l'actualité et pose le problème, des paragraphes, une conclusion. On répond à la question posée (\emph{should}, \emph{in che modo}, \emph{в какой мере}) avec des exemples précis et vérifiables, pris \textbf{hors de l'article}. Au thème, on traduit tout : une phrase laissée en blanc ne rapporte rien.

\textbf{Gestion du temps.} Le rapport 2026 rappelle qu'aucune durée n'est prescrite pour chaque exercice. Une répartition raisonnable : environ 10 minutes de lecture, 25 minutes pour le thème, 15 minutes pour la question 1, 35 minutes pour la question 2 et 5 minutes pour relire les accords et compter les mots. La question 2 pèse autant que le thème.

\textbf{Ce que dit le rapport 2026.} Le rapport de la session traite les six langues, de façon inégale : sections détaillées en allemand, anglais, arabe, espagnol et italien, une page en russe. Pour l'essai, le cadre commun est net : introduction qui contextualise et fait ressortir un problème, exemples probants, conclusion qui tire les leçons, et « La répétition de l'argumentaire présenté dans le sujet est déconseillée. » (p.~44). Les points que le jury donne pour \textbf{discriminants} sont :

\begin{itemize}
\tightlist
\item
  en anglais, les deux « depuis » du thème (\emph{past perfect} avec \emph{for}, \emph{present perfect} avec \emph{since}) et l'essai, jugé « particulièrement discriminant » par la richesse de la langue (p.~50 et 53) ;
\item
  en arabe, le respect du nombre de mots, dont le non-respect est « très sévèrement sanctionné » (p.~55) ; le jury de russe le relève aussi ;
\item
  en espagnol, la traduction du « vous » du thème (\emph{usted}), que l'accroche permettait de trancher ;
\item
  partout, les exemples personnels de l'essai, la réponse directe et non « catalogue » de la question 1.
\end{itemize}

Statistiques publiées : aucune moyenne, sauf en russe (notes de 10 à 16, moyenne 13,25). Le jury d'allemand estime qu'une copie complète et bien préparée dépasse 12/20, et celui d'italien juge la moyenne « très satisfaisante ».

\hypertarget{partie-i-allemand}{%
\section{Partie I --- Allemand}\label{partie-i-allemand}}

Article d'information sur les jeux vidéo à visée éducative. La question 1 en extrait les possibilités, la question 2 élargit à l'éducation aux valeurs. Le thème, au passé composé, mêle discours rapporté et dialogue.

\hypertarget{question-1-allemand-expression-uxe9crite-1}{%
\subsection{Question 1 --- Allemand, expression écrite 1)}\label{question-1-allemand-expression-uxe9crite-1}}

\begin{indication}{1}

La question porte sur les \emph{Möglichkeiten} : relevez dans l'article ce que ces jeux permettent, en laissant de côté ce qui relève du simple constat sur l'industrie.

\end{indication}

\begin{indication}{2}

L'article donne un mécanisme (l'interaction change le rapport au vécu), deux exemples nommés, un argument d'audience et une limite. Organisez ces quatre éléments en quatre ou cinq phrases.

\end{indication}

\begin{indication}{3}

Pour reformuler, changez la structure plutôt que les mots isolés : une relative généralisante (\emph{Wer \ldots{}}) ou une causale (\emph{weil}, \emph{da}) vous évite de recopier la citation de Samide.

\end{indication}

\hypertarget{solution}{%
\subsubsection{Solution}\label{solution}}

On garde quatre éléments : le mécanisme, les deux jeux (enquête sur les crimes nazis, lutte contre l'antisémitisme), l'audience (six personnes sur dix jouent) et la limite, exprimée par \emph{Allheilmittel}. On ne garde pas la statistique sur les livres, qui ne parle pas des jeunes.

Serious Games eröffnen der Bildungsarbeit neue Wege, weil sie Unterhaltung mit ernsten Themen verbinden. Wer selbst handelt, entwickelt eine engere Beziehung zum Erlebten als ein Leser oder Zuschauer. So lässt „The Darkest Files`` die Spieler NS-Verbrechen aufklären und zeigt, wie Unrecht ungestraft bleiben konnte; „Das Ilios-Experiment`` soll Schüler gegen Antisemitismus sensibilisieren und Toleranz fördern. Da sechs von zehn Deutschen spielen, erreicht das Medium ein sehr breites Publikum. Ein Wundermittel gegen rechtes Denken ist es zwar nicht, aber es kann Türen öffnen.

\textbf{(81 Wörter)}

\begin{piege}

N'ajoutez rien que le texte ne dise pas. Écrire que les jeunes ne lisent plus de livres, par exemple, déforme le propos : l'article dit seulement que beaucoup de gens passent plus de temps à jouer qu'à lire. En question 1, tout ajout de ce genre est compté comme une faute de compréhension.

\end{piege}

\begin{methode}

Rejet du verbe : après \emph{weil}, \emph{da} et \emph{wer}, le verbe conjugué va en fin de proposition. Quand la principale suit la subordonnée, elle commence par son verbe (\emph{Da sechs von zehn Deutschen spielen, \textbf{erreicht} das Medium \ldots{}}).

\end{methode}

\begin{jury}

Rapport 2026 (p.~47) : « il faut répondre à la question directement, sans introduction ni accroche », et noter le nombre de mots à la fin. La réponse proposée par le jury retient les mêmes éléments : transmission de savoirs et de valeurs humanistes, audience (six Allemands sur dix), immersion historique (« The Darkest Files »), antisémitisme et tolérance (« Das Ilios-Experiment »).

\end{jury}

\hypertarget{question-2-allemand-expression-uxe9crite-2}{%
\subsection{Question 2 --- Allemand, expression écrite 2)}\label{question-2-allemand-expression-uxe9crite-2}}

\begin{indication}{1}

La formule \emph{nicht nur \ldots, sondern auch} oriente vers le oui. La valeur du devoir tient donc à la manière de justifier ce oui et de traiter l'objection.

\end{indication}

\begin{indication}{2}

La consigne impose des exemples (\emph{anhand von Beispielen}). Préparez-en deux ou trois réels et extérieurs à l'article : un débat d'actualité allemand, un auteur dont vous savez tirer un argument, une pratique scolaire.

\end{indication}

\begin{indication}{3}

Pour présenter l'objection (« c'est le rôle de la famille ») sans la reprendre à votre compte, passez au discours indirect. Et gardez les deux termes du sujet : le savoir comme les valeurs.

\end{indication}

\hypertarget{solution-1}{%
\subsubsection{Solution}\label{solution-1}}

On suit ce plan : une introduction qui part de l'actualité (le choc PISA, la poussée de l'extrême droite) et pose la question, une thèse nuancée, deux arguments illustrés (Hans Jonas, l'apprentissage de la démocratie), l'objection rapportée au \emph{Konjunktiv I} et réfutée, puis une conclusion qui réunit savoir et valeurs.

Seit dem PISA-Schock von 2001 misst man deutsche Schulen vor allem an ihren Leistungen. Der Rechtsruck der letzten Jahre zeigt jedoch, dass Wissen allein die Demokratie nicht schützt. Muss die Schule also auch Werte vermitteln? Meiner Meinung nach ja, aber nicht auf Kosten des Wissens.

Erstens ist Wissen ohne Urteilskraft wirkungslos, ja sogar gefährlich. Ein Ingenieur muss rechnen können, aber auch die Folgen seiner Entscheidungen für Klima und Menschen abwägen. Das fordert Hans Jonas in „Das Prinzip Verantwortung``: Je mächtiger die Technik wird, desto größer wird unsere Verantwortung.

Zweitens muss Demokratie eingeübt werden. Viele Schulen tun das bereits: Klassensprecherwahlen, Besuche in KZ-Gedenkstätten, Gespräche mit Zeitzeugen. Da Bildung in Deutschland Ländersache ist, setzt jedes Land dabei eigene Schwerpunkte.

Man könnte zwar einwenden, die Wertevermittlung sei Sache der Familie. Doch jeder Unterricht vermittelt ohnehin Werte: durch die Auswahl der Texte, den Umgang mit Fehlern, die Regeln der Diskussion.

Werte lassen sich allerdings nicht predigen, sondern nur erfahren, und sie brauchen Wissen als Grundlage: Wer Nachhaltigkeit verteidigen will, muss zuerst den Klimawandel verstehen. Wissen und Werte gehören also zusammen: Erst beide machen aus Schülern mündige Bürger.

\textbf{(183 Wörter)}

\begin{methode}

Objection au discours indirect : \emph{Man könnte einwenden, die Wertevermittlung \textbf{sei} Sache der Familie}. Le \emph{Konjunktiv I} signale qu'on rapporte un avis sans le partager. La tournure ne coûte que quelques mots et montre une vraie maîtrise.

\end{methode}

\begin{piege}

\emph{Bildung} désigne l'éducation et la culture au sens large, \emph{Ausbildung} la formation professionnelle, \emph{Erziehung} l'éducation morale reçue en famille. Le sujet porte justement sur leur articulation : les confondre brouille tout le devoir.

\end{piege}

\begin{piege}

Fautes relevées par le jury : le genre (\emph{der Unterschied}, \emph{der Einfluss}, \emph{die Rolle}, \emph{das Bewusstsein}, \emph{das Wissen}), \emph{Werte als Demokratie} au lieu de \emph{Werte wie Demokratie}, \emph{die Jungen} ou \emph{die junge Menschen} au lieu de \emph{junge Menschen}, \emph{die Jugendlichen}, \emph{die Jugend}. On ne dit pas non plus \emph{Werte lehren} mais \emph{Werte vermitteln}.

\end{piege}

\begin{jury}

Rapport 2026 (p.~47) : une accroche d'actualité est « toujours un bonus », mais « attention ensuite à ne pas renverser la question en ne traitant qu'une partie du sujet (le savoir ou les valeurs) ». Un auteur cité doit servir un argument, et la connaissance du système éducatif allemand, qui dépend des Länder, est bienvenue. Le jury regrette enfin qu'il n'y ait « pas assez de marquage du positionnement personnel » : \emph{Meiner Meinung nach}, \emph{Meines Erachtens}, \emph{Für mich} + verbe.

\end{jury}

\hypertarget{question-3-allemand-thuxe8me}{%
\subsection{Question 3 --- Allemand, thème}\label{question-3-allemand-thuxe8me}}

\begin{indication}{1}

Repérez les narrateurs : c'est une femme qui parle (\emph{venue}, \emph{allées}), et la tante prend la parole dans la première phrase. Repérez aussi ce qui est cité directement et ce qui est rapporté.

\end{indication}

\begin{indication}{2}

Dans la réplique de la tante, \emph{elle dit qu'elle ne veut \ldots{} que ça ne sert à rien} est du discours rapporté : pensez au mode qui convient en allemand écrit.

\end{indication}

\begin{indication}{3}

Les difficultés lexicales sont \emph{sermonner}, \emph{ça ne sert à rien}, \emph{j'ai laissé tomber} et \emph{ce n'était pas mon affaire}. Cherchez à chaque fois l'expression idiomatique allemande plutôt que le calque.

\end{indication}

\hypertarget{solution-2}{%
\subsubsection{Solution}\label{solution-2}}

„Du, die du Lebenserfahrung hast``, sagte meine Tante an mich gewandt -- von welcher Erfahrung sprach sie? --, „du solltest ihr sagen, sie solle sich bessern. Sie tut nichts mehr für die Schule, sie sagt, sie wolle nicht einmal weitermachen, das bringe ja doch nichts.`` {[}\ldots{]}

Eines Nachmittags habe ich ihr beim Schulausgang ins Gewissen geredet; ich war gekommen, um dort auf sie zu warten. Wir sind in ein Café gegangen {[}\ldots{]}. Sie hat sich auf die Terrasse gesetzt, um rauchen zu können.

„Vielleicht solltest du nicht schon so jung rauchen``, habe ich zu ihr gesagt.

„Weil du etwa nicht rauchst?{}``

„In deinem Alter rauchte ich noch nicht.``

„Und welchen Unterschied macht das jetzt?{}``

Ich habe es dabei bewenden lassen. Schließlich ging es mich nichts an, ob sie in der Öffentlichkeit oder heimlich rauchte.

Choix justifiés :

\begin{itemize}
\tightlist
\item
  \emph{Toi qui as \ldots{}} devient \emph{Du, die du \ldots{} hast} : on reprend le pronom \emph{du} dans la relative, sinon l'accord du verbe est faux.
\item
  Le discours rapporté passe au \emph{Konjunktiv I} : \emph{sie wolle}, \emph{das bringe}.
\item
  \emph{sermonner} se rend par \emph{jemandem ins Gewissen reden} ; \emph{j'ai laissé tomber} par \emph{es dabei bewenden lassen}, dont le double infinitif au parfait est \emph{bewenden lassen} et non \emph{gelassen}.
\item
  \emph{ce n'était pas mon affaire} se rend par \emph{es ging mich nichts an}.
\item
  \emph{elle ne travaille plus à l'école} devient \emph{sie tut nichts mehr für die Schule} : \emph{sie arbeitet nicht mehr in der Schule} ferait croire qu'elle y est employée.
\item
  \emph{Parce que toi tu ne fumes pas ?} garde la subordonnée : \emph{Weil du etwa nicht rauchst?}, verbe en fin de proposition, comme dans le corrigé du jury.
\item
  \emph{un après-midi} est un repère ponctuel, donc un génitif : \emph{eines Nachmittags} (\emph{einen Nachmittag} dirait la durée). \emph{à la sortie de l'école} se rend par \emph{beim Schulausgang}.
\item
  Le récit est au passé composé : \emph{habe ich ihr ins Gewissen geredet}, \emph{habe ich zu ihr gesagt}. Dans le dialogue, \emph{je ne fumais pas} passe au prétérit, \emph{rauchte}.
\end{itemize}

\begin{piege}

\emph{korrigieren} s'emploie pour une copie, pas pour une conduite : \emph{sich bessern}. \emph{das dient zu nichts} est un calque ; on écrit \emph{das bringt nichts} ou \emph{das nützt nichts}.

\end{piege}

\begin{piege}

Les fautes de forme que le jury relève le plus : \emph{sprachte} pour \emph{sprach}, \emph{lässte} pour \emph{ließ}, \emph{rauchtete} pour \emph{rauchte}, \emph{gesetz} pour \emph{gesetzt}, \emph{gekommt} pour \emph{gekommen}. Viennent ensuite \emph{Du solltest zu ihr zu sagen} (pas de \emph{zu} après un verbe de modalité), \emph{für sie warten} au lieu de \emph{auf sie warten}, et \emph{an die Terrasse} au lieu de \emph{auf die Terrasse}.

\end{piege}

\begin{jury}

Rapport 2026 (p.~45-46) : le jury classe les fautes « des plus lourdes aux plus acceptables » et commence par la place du verbe conjugué, puis les formes verbales. Sa proposition de traduction rend \emph{j'ai laissé tomber} par \emph{Ich gab dann auf} et \emph{sermonner} par \emph{tadeln} ; \emph{ins Gewissen reden} et \emph{es dabei bewenden lassen} sont des variantes idiomatiques.

\end{jury}

\hypertarget{partie-ii-anglais}{%
\section{Partie II --- Anglais}\label{partie-ii-anglais}}

L'article analyse la perte de compétitivité d'un distributeur en ligne. La question 1 impose une reformulation, la question 2 un débat sur la consommation. Le thème, un récit au passé, est un exercice sur l'emploi des temps.

\hypertarget{question-4-anglais-expression-uxe9crite-1.}{%
\subsection{Question 4 --- Anglais, expression écrite 1.}\label{question-4-anglais-expression-uxe9crite-1.}}

\begin{indication}{1}

La consigne dit \emph{in your own words} : prévoyez des équivalents pour \emph{lumbered with a stock hangover}, \emph{undercut} et \emph{nimble}.

\end{indication}

\begin{indication}{2}

Classez les causes plutôt que de les énumérer : une erreur stratégique passée, des clients qui dépensent autrement, une concurrence plus agile (prix, technologie, occasion), puis une incapacité à investir.

\end{indication}

\begin{indication}{3}

Faites attention aux temps. La cause est passée (\emph{misjudged}) et ses conséquences sont présentes : le \emph{present perfect} et le présent servent à faire le lien.

\end{indication}

\hypertarget{solution-3}{%
\subsubsection{Solution}\label{solution-3}}

Asos misjudged the post-Covid market: expecting lockdown shoppers to stay online for good, it overstocked and was left with £1bn of unsold clothes once people returned to shops. Customers now also have less to spend and prefer holidays, festivals or streaming to new outfits. Meanwhile, rivals innovate faster: Shein, helped by a tax exemption, sells for less and uses technology to spot trends quickly, while Vinted attracts eco-conscious buyers. Squeezed by falling sales, higher wages and new regulations, Asos cannot afford to keep up.

\textbf{(84 words)}

\begin{piege}

\emph{tax break} désigne un avantage fiscal légal. Parler de \emph{loophole} ou de fraude, c'est faire dire au texte plus qu'il ne dit. Ne remplacez pas non plus \emph{£1bn} par \emph{one milliard}.

\end{piege}

\begin{piege}

Le jury dit avoir cessé de compter les réponses qui commencent par \emph{According to the Guardian} ou qui recopient la question (\emph{Asos is currently struggling to remain competitive because\ldots{}}), ainsi que les réponses en plusieurs paragraphes avec introduction. Entrez directement dans les causes.

\end{piege}

\begin{jury}

Rapport 2026 (p.~51-52) : « Le jury attend une réponse directe, concise et précise à la question. » Il faut repérer quatre ou cinq éléments : l'après-Covid mal anticipé et le stock d'invendus, des concurrents (Shein, Vinted) plus innovants et moins chers, l'exemption de taxes de Shein, des clients qui dépensent autrement et regardent le prix ou la durabilité. Reformuler ne veut pas dire commenter : « Il s'agit de retranscrire les faits et uniquement les faits. »

\end{jury}

\begin{methode}

Pour reformuler, on change la catégorie grammaticale (\emph{had cash to spare} devient \emph{could afford}) ou le point de vue (\emph{Shein undercut Asos} devient \emph{Shein sells for less}). Remplacer un mot par son synonyme, sans rien changer d'autre, ne suffit pas.

\end{methode}

\hypertarget{question-5-anglais-expression-uxe9crite-2.}{%
\subsection{Question 5 --- Anglais, expression écrite 2.}\label{question-5-anglais-expression-uxe9crite-2.}}

\begin{indication}{1}

Le mot important est \emph{only}, et la question est au modal \emph{should} : répondez à ce modal, pas à la question « le prix est-il le seul critère ? ».

\end{indication}

\begin{indication}{2}

La concession doit être sérieuse : pour un ménage pauvre, le prix est une contrainte et non une préférence. Le contexte d'inflation et de crise du pouvoir d'achat dans le monde anglophone fournit l'introduction.

\end{indication}

\begin{indication}{3}

Cherchez l'argument central dans l'idée de coût déplacé : qui paie le prix bas, et à quel moment ? Illustrez-le par un fait et une référence culturelle pris hors de l'article, et laissez Shein et Vinted de côté.

\end{indication}

\hypertarget{solution-4}{%
\subsubsection{Solution}\label{solution-4}}

Plan : une introduction qui part de l'actualité (inflation, élection de Zohran Mamdani à New York sur le pouvoir d'achat) et pose le problème, une prise de position, une concession, l'argument du coût déplacé (Rana Plaza, \emph{Brave New World}), puis une conclusion qui répond.

With inflation squeezing households on both sides of the Atlantic, affordability has become a political issue: Zohran Mamdani won New York's mayoral race in 2025 on that very promise. Can consumers still afford to look beyond the price tag, and should they? I believe they should, even if price comes first for many.

Admittedly, for the poorest families price is a constraint, not a choice. When rent and energy bills swallow most of an income, lecturing them about ethics is unfair: only the well-off can pay more.

Yet a low price is often a cost shifted onto someone else. The collapse of the Rana Plaza factory in Bangladesh, which killed more than 1,100 garment workers in 2013, revealed who really pays for bargain clothes. Cheap goods that fall apart quickly also cost more in the long run, as Huxley foresaw in \emph{Brave New World}, whose citizens are taught that ``ending is better than mending''.

Consumers should therefore weigh durability, working conditions and environmental impact too, and governments should make these choices affordable to all. Price matters, but it should never be the only yardstick.

\textbf{(184 words)}

\begin{piege}

Le pluriel de \emph{criterion} est \emph{criteria} ; \emph{criterias} n'existe pas. \emph{Actually} veut dire « en fait » et \emph{eventually} « finalement ».

\end{piege}

\begin{piege}

Reprendre les exemples de l'article. Le jury constate que trop de candidats ont repris Shein et Vinted, ce qui rendait le propos « plat et convenu ». Les meilleures copies sont allées vers l'éthique (Musk et Tesla, Bezos et Amazon), vers l'actualité (inflation, crise du logement au Royaume-Uni) ou vers des références culturelles (Huxley, \emph{Le Diable s'habille en Prada}).

\end{piege}

\begin{jury}

Rapport 2026 (p.~52-53) : courte introduction qui contextualise et problématise, deux ou trois paragraphes avec une idée directrice et des exemples concrets, courte conclusion. « Une prise de position claire et nuancée est un critère essentiel. » Les exemples doivent dépasser le texte : « la reprise d'exemples de l'article est à proscrire ». L'exercice est jugé « particulièrement discriminant » par la richesse de la langue.

\end{jury}

\hypertarget{question-6-anglais-thuxe8me}{%
\subsection{Question 6 --- Anglais, thème}\label{question-6-anglais-thuxe8me}}

\begin{indication}{1}

Commencez par fixer le repère temporel du récit : \emph{venait alors d'avoir seize ans}. Ce qui le précède se met au \emph{past perfect}, ce qui se situe à ce moment au \emph{simple past}.

\end{indication}

\begin{indication}{2}

L'imparfait de répétition (\emph{ne cessait de}, \emph{répétait-elle}) demande une forme fréquentative. Dans le dialogue, \emph{Je te connais depuis que \ldots{}} appelle un temps différent du français.

\end{indication}

\begin{indication}{3}

Pour \emph{Elle voulait qu'on la laisse tranquille}, l'anglais n'accepte pas \emph{wanted that} : passez par un passif ou par une proposition infinitive.

\end{indication}

\hypertarget{solution-5}{%
\subsubsection{Solution}\label{solution-5}}

Camille had just turned sixteen at the time. For several weeks, her behaviour had been worrying her parents. Although she had always been a brilliant pupil, she was now reluctant to go to school. Her mother kept asking her whether something had happened. ``No, nothing, no, nothing,'' she would repeat. She no longer wanted to get up. One day, she had finally murmured: ``I feel a weight inside me that is impossible to lift.'' Her mother had already heard that kind of remark at the hospital, from patients suffering from depression. Isabelle assumed that she needed to help her daughter on a daily basis; that way, the girl would get her energy back. {[}\ldots{]}

Sitting next to Camille, the doctor did not know what to say.

``You can tell me anything. You know very well that I'm a friend of the family. I've known you since you were a little girl.''

``I know.''

``Then tell me what's the matter. Tell me where it hurts.''

``It doesn't hurt,'' said Camille. {[}\ldots{]}

She wanted to be left alone.

Choix justifiés :

\begin{itemize}
\tightlist
\item
  \emph{depuis plusieurs semaines \ldots{} inquiétait} se rend par \emph{had been worrying} : la durée se situe avant le repère.
\item
  \emph{ne cessait de demander} devient \emph{kept asking} et \emph{répétait-elle} devient \emph{she would repeat}.
\item
  \emph{s'imagina} est un passé simple, donc un \emph{simple past} (\emph{assumed}), et non \emph{had assumed}.
\item
  \emph{Je te connais depuis que} demande le \emph{present perfect} : \emph{I've known you since}.
\item
  \emph{Je n'ai pas mal} devient \emph{It doesn't hurt}, en écho à \emph{where it hurts}.
\item
  Le texte source omet un pronom (\emph{Tu sais bien que suis un ami}) ; on traduit la phrase attendue.
\item
  \emph{Elle qui avait toujours été \ldots{} rechignait} oppose ce que Camille a été et ce qu'elle est devenue : une concessive, \emph{Although she had always been \ldots, she was now reluctant to}, rend ce contraste.
\item
  \emph{rechigner} se rend par \emph{to be reluctant to} (\emph{to drag her feet about} + V-ing est encore plus idiomatique) ; \emph{to refuse} sous-traduit. \emph{brillante} : \emph{brilliant} convient, \emph{outstanding} ou \emph{a straight-A student} sont valorisés.
\item
  \emph{toute petite} parle de l'âge : \emph{little} (\emph{a little girl}), et non \emph{small}, qui parle de la taille.
\end{itemize}

\begin{methode}

Avant de traduire un récit, tracez une frise : un repère, puis l'antériorité (\emph{past perfect}), la simultanéité (\emph{simple past} ou \emph{past continuous}) et la répétition (\emph{would}, \emph{kept} + V-ing).

\end{methode}

\begin{piege}

\emph{fit Camille} signifie « dit Camille » : \emph{said}, \emph{answered}, jamais \emph{did}. Le dialogue anglais n'a ni tiret ni chevrons : chaque réplique s'ouvre et se ferme par ``\ldots{}''. Le passer au discours indirect n'a pas de sens en traduction. Enfin, \emph{she didn't wanted} ou \emph{she didn't knew} sont des fautes de base que le jury sanctionne sans indulgence.

\end{piege}

\begin{jury}

Rapport 2026 (p.~49-51) : la différence s'est faite d'entrée sur \emph{venir d'avoir seize ans}, \emph{to turn sixteen}. Sur les deux « depuis » (\emph{had been worrying \ldots{} for several weeks}, \emph{I've known you since}), le jury note que « ces deux points de grammaire furent discriminants ». Il attendait \emph{she wanted to be left alone} ou \emph{she wanted everyone to leave her alone}.

\end{jury}

\hypertarget{partie-iii-arabe}{%
\section{Partie III --- Arabe}\label{partie-iii-arabe}}

Article d'opinion sur le mouvement de la jeunesse marocaine en 2025. On répond en arabe littéral (فصحى), avec ses propres mots. Le thème doit être écrit sans voyelles et oblige à maîtriser le duel.

\hypertarget{question-7-arabe-expression-uxe9crite-ux627ux644ux633ux624ux627ux644-ux627ux644ux623ux648ux644}{%
\subsection{Question 7 --- Arabe, expression écrite, السؤال الأول}\label{question-7-arabe-expression-uxe9crite-ux627ux644ux633ux624ux627ux644-ux627ux644ux623ux648ux644}}

\begin{indication}{1}

Le texte présente les causes de façon dispersée, en paragraphes 2 à 5. Repérez-les : crise des élites, perte de crédit des corps intermédiaires, information circulant en ligne, coût de la vie, inégalités, nouvelle culture politique.

\end{indication}

\begin{indication}{2}

La consigne exige vos propres expressions (مستعملا عباراتك الخاصة). Remplacez par exemple « عجزها عن التكيف » par un synonyme et changez la construction des phrases.

\end{indication}

\begin{indication}{3}

Regroupez les causes en deux temps, d'abord les causes politiques puis les causes sociales et numériques, et reliez-les par des connecteurs (أولا، ثانيا، يضاف إلى ذلك).

\end{indication}

\hypertarget{solution-6}{%
\subsubsection{Solution}\label{solution-6}}

يرى الكاتب أن احتجاجات «جيل زد» في المغرب اندلعت لأسباب متشابكة. فهي أولا نتيجة لأزمة النخب السياسية التقليدية العاجزة عن مواكبة التحولات، ولضعف الأحزاب والنقابات التي فقدت دورها في الوساطة، فلم يعد الشباب يثقون بخطاب يرونه مكررا لا يوفر لهم عملا ولا خدمات عامة لائقة. وثانيا، كشفت لهم شبكة الإنترنت ما يتمتع به المسؤولون من امتيازات وتورط بعضهم في الفساد، في وقت يعانون فيه من غلاء المعيشة. يضاف إلى ذلك تفاوت اجتماعي صارخ بين الجهات، وثقافة سياسية جديدة اكتسبوها من الشبكات الاجتماعية.

\textbf{(82 كلمة)}

\begin{piege}

Un verbe placé avant son sujet reste au singulier : تراجعت المؤسسات, et non تراجعن. Évitez aussi les pronoms redondants comme « يواجهون هم ». Pour compter, on prend les mots graphiques : و، ف، ب، ل et ال, soudés au mot, ne comptent pas à part (le rapport 2026 le rappelle pour و).

\end{piege}

\begin{jury}

Rapport 2026 (p.~54) : la réponse se fonde uniquement sur le texte ; « une réponse exhaustive n'est pas exigée mais le jury attend que les candidats dégagent au moins les causes principales évoquées par le texte ». Le jury a sanctionné, parfois lourdement, les réponses « catalogue » et la reprise des expressions du texte sans reformulation. Pas d'introduction ni de conclusion obligatoires, et aucun jugement personnel ni information extérieure.

\end{jury}

\hypertarget{question-8-arabe-expression-uxe9crite-ux627ux644ux633ux624ux627ux644-ux627ux644ux62bux627ux646ux64a}{%
\subsection{Question 8 --- Arabe, expression écrite, السؤال الثاني}\label{question-8-arabe-expression-uxe9crite-ux627ux644ux633ux624ux627ux644-ux627ux644ux62bux627ux646ux64a}}

\begin{indication}{1}

Le sujet propose une alternative : il vaut mieux déplacer le critère que choisir l'un des deux termes sans nuance. Demandez-vous ce qui est vraiment nouveau, l'intensité de la contestation ou la nature des revendications.

\end{indication}

\begin{indication}{2}

Les exemples doivent venir de la réalité arabe et être datés : soulèvements de 2011, mouvements de 2016 à 2019 au Maroc, au Liban et en Irak, mobilisation marocaine de 2025.

\end{indication}

\begin{indication}{3}

Ouvrez par une ou deux phrases qui situent les manifestations et reformulent l'alternative. La conclusion peut lier l'avenir du mouvement à la question que le texte pose en dernier : la réforme des partis.

\end{indication}

\hypertarget{solution-7}{%
\subsubsection{Solution}\label{solution-7}}

منذ نهاية سبتمبر 2025، يتظاهر شباب مغاربة بدعوة من مجموعات رقمية. فهل هو تحول جذري في علاقتهم بالسياسة أم موجة عابرة؟ في رأيي، تمثل حركة «جيل زد» تحولا حقيقيا في الوعي السياسي، لكنه تحول هش لم يكتمل بعد.

فهي أولا تحول حقيقي لأنها غيرت طبيعة المطالب نفسها. فهذا الجيل لا ينطلق من إيديولوجيا كبرى، بل من المستشفى والمدرسة وفرص العمل. وقد ظهر ذلك حين قارن الشباب بين الأموال المخصصة للملاعب استعدادا لكأس العالم وبين حالة المستشفيات العمومية. وهذا هو الفرق بين 2011 واليوم: فشعار «الشعب يريد إسقاط النظام» كان شاملا، أما مطالب «جيل زد» فملموسة وقابلة للقياس. كما أن وسائل التواصل جعلت الشباب العربي يتابع ما يجري في البلدان الأخرى لحظة بلحظة، ويستلهم منها الشعارات وأساليب التنظيم.

ثانيا، تدعو تجارب المنطقة إلى الحذر. فحراك الريف في المغرب سنة 2016 انتهى باعتقال قادته، والحراك اللبناني سنة 2019 لم يغير الطبقة السياسية، واحتجاجات تشرين في العراق لم تحقق إلا جزءا من مطالبها. فالشبكات الاجتماعية تحشد بسرعة، لكنها لا تفرز قيادة ولا برنامجا.

لذلك أعتقد أن مستقبل هذه الحركة مرهون بقدرتها على التنظيم، وبقدرة الأحزاب على الإصلاح: فإذا تجددت المؤسسات تحول الغضب إلى سياسة عامة، وإلا عاد الشباب إلى الشارع مرة أخرى.

\textbf{(189 كلمة)}

\begin{methode}

Pour une question du type « mutation durable ou phénomène passager ? », on accorde à chaque terme ce qui lui revient, en distinguant la forme, nouvelle, et la durée, incertaine. On tranche à la fin par une condition (إذا \ldots{} وإلا).

\end{methode}

\begin{piege}

N'inventez ni chiffre de participation ni slogan précis. Si vous n'êtes pas sûr d'un fait, restez sur des faits établis : arrestation des meneurs du Hirak du Rif, démission du gouvernement irakien après Tichrine. Évitez aussi les tournures dialectales (ما كاين، ديال).

\end{piege}

\begin{piege}

Deux hors-sujet relevés par le jury : le réquisitoire contre le collectif (accusé de troubler l'ordre public ou d'être manipulé de l'étranger), qui ne répond pas à la question ; et l'évocation du printemps arabe sans dire ce qui relève de la continuité et ce qui relève de la rupture. Reproduire les exemples de l'article sans rien y ajouter a aussi été fermement sanctionné.

\end{piege}

\begin{jury}

Rapport 2026 (p.~54-55) : introduction même très brève qui contextualise et reformule, développement argumenté, conclusion. « Le jury se désole de l'augmentation significative du nombre de copies dont la réponse est constituée d'une seule traite sans paragraphe » et les a lourdement pénalisées. Il attend une approche nuancée « indépendamment du point de vue défendu », et rappelle que le non-respect du nombre de mots est « très sévèrement sanctionné ».

\end{jury}

\hypertarget{question-9-arabe-thuxe8me}{%
\subsection{Question 9 --- Arabe, thème}\label{question-9-arabe-thuxe8me}}

\begin{indication}{1}

Comptez les acteurs : « mes parents » sont deux. Le pronom, le verbe et l'adjectif doivent être au duel partout où ils les reprennent.

\end{indication}

\begin{indication}{2}

Pour \emph{ils nous l'assénèrent}, choisissez une forme qui dise la répétition et l'insistance. La suite nominale \emph{Ongles propres, dents brossées\ldots{}} peut rester une suite de groupes nominaux.

\end{indication}

\begin{indication}{3}

Gardez la métaphore militaire de \emph{partir en mission}, l'anaphore \emph{aucune \ldots{} aucun \ldots{} aucune} et le ton de \emph{tant pis}.

\end{indication}

\hypertarget{solution-8}{%
\subsubsection{Solution}\label{solution-8}}

منذ أن أقنعتهما مديرة روضتنا بأن الجمع بين لغتين سيقودنا إلى الفشل الدراسي، أقصى والداي اللغة العربية من البيت. فالفشل الدراسي في نظرهما أعظم مأساة على الإطلاق. أما التفوق في المدرسة، وقد ظلا يرددانه على مسامعنا طوال سنوات دراستنا، فهو أكثر بكثير من ضرورة: إنه واجب علينا تجاه شهداء الثورة الجزائرية. أظافر نظيفة، وأسنان منظفة بالفرشاة، وقمصان مكوية، وشعر ممشط: يحرص والداي على أن نظهر للفرنسيين أفضل صورة عن أنفسنا. لن يقبلا أي حماقة، ولا أي ملاحظة في دفتر المراسلة، ولا أي علامة سيئة. نحن لا نذهب إلى المدرسة فحسب، بل نمضي في مهمة. ولا يهم إن كان لا بد من التضحية بالعربية. وهكذا أرسيت عادة ثابتة: قراءة قصة عند النوم، وإملاء صباح السبت، واستظهار قصيدة بعد ظهر الأحد.

Choix justifiés :

\begin{itemize}
\tightlist
\item
  Le duel est employé partout : أقنعتهما، في نظرهما، ظلا يرددانه، لن يقبلا. Le verbe antéposé reste au singulier : أقصى والداي.
\item
  \emph{école maternelle} se rend par روضة, qui désigne le jardin d'enfants.
\item
  \emph{bannir} se rend par أقصى, sans la nuance religieuse que prendrait حرم.
\item
  La négation absolue est rendue par لن يقبلا أي \ldots{} ولا أي \ldots{} ولا أي.
\item
  \emph{Un rituel} se rend par عادة, comme le propose le jury ; عبادة (culte religieux) est un non-sens.
\item
  Le texte n'est pas vocalisé, conformément à l'astérisque de l'énoncé.
\end{itemize}

\begin{piege}

Vocaliser une partie du texte est plus risqué que ne rien vocaliser, puisque chaque voyelle écrite peut alors être vérifiée. L'énoncé interdit d'ailleurs la vocalisation.

\end{piege}

\begin{piege}

Les contresens lexicaux relevés par le jury : الضحايا pour « les martyrs » (الشهداء), النهضة pour « la révolution » (الثورة), فساتين pour « chemises » (أقمصة, ou قمصان, autre pluriel de قميص), بعثة pour « mission » (مهمة), عرض pour « récitation » (استظهار), et قراءة تاريخ, calque de « lecture d'une histoire », au lieu de قراءة قصة ou حكاية.

\end{piege}

\begin{jury}

Rapport 2026 (p.~55) : beaucoup de traductions partielles ou littérales tiennent à une maîtrise insuffisante du français, et « toute omission est fortement sanctionnée ». Les meilleures copies respectent les particularités des deux langues.

\end{jury}

\hypertarget{partie-iv-espagnol}{%
\section{Partie IV --- Espagnol}\label{partie-iv-espagnol}}

Article chiffré sur l'échec scolaire des jeunes Gitans. La question 2 est la seule des six à demander des connaissances : elle porte sur les mesures prises des deux côtés de l'Atlantique. Le thème, au présent, porte sur le subjonctif et sur l'opposition entre \emph{ser} et \emph{estar}.

\hypertarget{question-10-espagnol-expression-uxe9crite-1.}{%
\subsection{Question 10 --- Espagnol, expression écrite 1.}\label{question-10-espagnol-expression-uxe9crite-1.}}

\begin{indication}{1}

Ne transformez pas la réponse en liste de pourcentages. Cherchez l'enchaînement des obstacles, de la petite enfance jusqu'à l'université.

\end{indication}

\begin{indication}{2}

Le texte cite aussi des obstacles qui ne sont pas chiffrés : regroupez-les en une phrase, reformulée. Terminez par la conséquence, pour l'université et pour l'avenir professionnel.

\end{indication}

\hypertarget{solution-9}{%
\subsubsection{Solution}\label{solution-9}}

Los gitanos sufren una brecha educativa que empieza en la infancia: pocos niños asisten a la educación infantil, y lo hacen tarde. Las repeticiones en primaria provocan luego un fuerte desfase curricular, y la mayoría no termina la ESO: el fracaso escolar alcanza el 62,8 \%. A ello se suman la segregación escolar, los prejuicios, las bajas expectativas y la falta de referentes gitanos con estudios, que dificultan su integración. Así, solo un 0,8 \% obtiene un título universitario, lo que compromete su futuro profesional.

\textbf{(83 palabras)}

\begin{piege}

En espagnol, les noms de peuples s'écrivent avec une minuscule (\emph{los gitanos}), et le pourcentage prend un article (\emph{el 63 \%}, \emph{un 0,8 \%}). Le verbe \emph{enfrentarse} se construit avec \emph{a} ou \emph{con}.

\end{piege}

\begin{piege}

Pas de raccourci. Le jury a lu que « los profesores y los alumnos se oponen a la presencia de gitanos en las escuelas » : l'article ne dit rien de tel, il rapporte une anecdote et des préjugés. Évitez aussi la liste sans connecteurs, qui donne une réponse « catalogue ».

\end{piege}

\begin{jury}

Rapport 2026 (p.~56) : une réponse complète met en avant cinq idées, « el fracaso escolar importante », « el retraso en el currículo », « la falta de modelos », « la dificultad de integración » et « el futuro profesional incierto ». La contextualisation n'est pas indispensable, et « Le manque de connecteurs conduit souvent à une restitution peu claire de l'information ».

\end{jury}

\hypertarget{question-11-espagnol-expression-uxe9crite-2.}{%
\subsection{Question 11 --- Espagnol, expression écrite 2.}\label{question-11-espagnol-expression-uxe9crite-2.}}

\begin{indication}{1}

\emph{Explique las medidas} demande un exposé documenté, pas une opinion sur la tolérance. Un plan géographique, l'Espagne puis l'Amérique latine, s'impose.

\end{indication}

\begin{indication}{2}

Pour l'Espagne, pensez à ce qui suit l'élève, à la stratégie de l'État, puis aux autres minorités (langues co-officielles, migrants). Pour l'Amérique latine, pensez aux peuples indigènes, à leurs langues et aux constitutions.

\end{indication}

\begin{indication}{3}

N'avancez que des dispositifs que vous savez nommer avec certitude. Dans une question de connaissances, deux exemples sûrs valent mieux que cinq approximatifs.

\end{indication}

\hypertarget{solution-10}{%
\subsubsection{Solution}\label{solution-10}}

Según el artículo, solo el 0,8 \% de los jóvenes gitanos obtiene un título universitario. Integrar a las minorías sigue siendo, pues, un reto a ambos lados del Atlántico, que se ha abordado por dos vías: el acompañamiento y el reconocimiento.

En España, la Fundación Secretariado Gitano desarrolla el programa Promociona, que orienta a los alumnos gitanos y a sus familias para que terminen la ESO y sigan estudiando, y el Estado aprobó una Estrategia Nacional para la Igualdad, Inclusión y Participación del Pueblo Gitano 2021-2030. Más allá de los gitanos, desde 2023 se pueden usar las lenguas cooficiales en el Congreso, y el Gobierno de Pedro Sánchez ha impulsado la regularización extraordinaria de migrantes.

En América Latina, donde las minorías son sobre todo pueblos indígenas, el eje ha sido el reconocimiento. La Constitución boliviana de 2009 define un Estado plurinacional y reconoce como oficiales decenas de lenguas indígenas, mientras que México ha creado universidades interculturales y Perú desarrolla la educación intercultural bilingüe.

Estas medidas muestran que la inclusión exige a la vez recursos concretos en las aulas y el reconocimiento de la identidad de cada minoría.

\textbf{(186 palabras)}

\begin{methode}

\emph{para que} est toujours suivi du subjonctif (\emph{para que terminen}). Dans une réponse d'exposé, c'est la construction subordonnée qui se place le plus naturellement et qui rapporte le plus.

\end{methode}

\begin{piege}

Dénoncer la situation des minorités ne répond pas à la question : le jury reproche à de nombreuses copies d'avoir donné des exemples de discrimination sans mesure concrète d'inclusion. Il épingle aussi des arguments hors de propos (budget d'un club de football, narcotrafic) et des fautes lexicales comme \emph{indígena} appliqué à Ibiza.

\end{piege}

\begin{jury}

Rapport 2026 (p.~56-57) : « Le jury s'attendait à ce que les candidats élargissent la problématique et ne se cantonnent pas à la communauté gitane. » On attend une idée, un exemple, une brève analyse, reliés par des connecteurs. Le jury cite comme pertinents le rappeur gitan Omar Montes et Bad Bunny au Super Bowl 2026 (la culture comme vecteur d'intégration), l'ONCE, l'officialisation de langues indigènes en Amérique latine, le projet de régularisation des migrants du gouvernement Sánchez et l'usage des langues co-officielles au Parlement.

\end{jury}

\hypertarget{question-12-espagnol-thuxe8me}{%
\subsection{Question 12 --- Espagnol, thème}\label{question-12-espagnol-thuxe8me}}

\begin{indication}{1}

Ne traduisez pas la phrase de contextualisation en italique (note de l'énoncé), mais lisez-la : elle dit à qui s'adresse le narrateur, donc qui est désigné par \emph{vous}.

\end{indication}

\begin{indication}{2}

Deux tournures du texte imposent le subjonctif imparfait. Repérez-les, ainsi que les phrases où il faut choisir entre \emph{ser} et \emph{estar}.

\end{indication}

\begin{indication}{3}

\emph{Il y en a qui \ldots{}} et \emph{on voit la peur dans leurs yeux} ont des équivalents idiomatiques, et la maison marocaine a un toit-terrasse.

\end{indication}

\hypertarget{solution-11}{%
\subsubsection{Solution}\label{solution-11}}

En la mina, mis compañeros y yo cogemos una máquina que nos hace descender bajo tierra. A menudo bajo a casi mil metros. En el vientre de la tierra. Los hay que vomitan, otros que tienen miedo, y se les ve el miedo en los ojos. A veces yo también tengo miedo. Pero cuando pienso en usted y rezo, mi corazón se serena gracias a Dios. Eso me da fuerzas. Nos hundimos en la oscuridad, es como si la tierra nos llamara. Es un viaje. Es un mundo extraño. Estoy fascinado y, al mismo tiempo, a menudo estoy preocupado. No tanto por mí sino por los demás. Nunca habría imaginado que bajo nuestros pies hubiera todo un mundo. Como cuando era niño y dormía en la azotea de la casa para descubrir la inmensidad del cielo.

Choix justifiés :

\begin{itemize}
\tightlist
\item
  \emph{comme si \ldots{} appelait} se rend par \emph{como si \ldots{} llamara} : après \emph{como si}, on emploie toujours le subjonctif imparfait.
\item
  Pour \emph{Je n'aurais jamais imaginé que \ldots{} il y avait}, on écrit \emph{nunca habría imaginado que \ldots{} hubiera}. L'indicatif \emph{había} se rencontre aussi, mais le subjonctif est la solution la plus sûre.
\item
  Il faut \emph{Es un viaje} (\emph{ser} devant un nom) mais \emph{Estoy fascinado} et \emph{estoy preocupado} (\emph{estar} pour un état).
\item
  \emph{vous} : l'accroche dit que le narrateur envoie des enregistrements \textbf{à son père}. C'est un vouvoiement de respect adressé à une seule personne : \emph{usted}, d'où \emph{cuando pienso en usted y rezo} (\emph{pensar en}, et deux subordonnées simplement coordonnées par \emph{y}).
\item
  \emph{Pas tellement pour moi, mais pour les autres} devient \emph{No tanto por mí sino por los demás}, la traduction attendue par le jury : après une négation, « mais » se rend par \emph{sino}. \emph{preocupado por} : \emph{por} suit les adjectifs de sentiment.
\item
  \emph{le toit} d'une maison marocaine est une \emph{azotea} ; \emph{tejado} désigne un toit de tuiles en pente.
\item
  \emph{on voit la peur dans leurs yeux} : \emph{se les ve el miedo en los ojos} ; le jury accepte aussi \emph{uno puede ver el miedo}.
\end{itemize}

\begin{piege}

Traduire \emph{vous} par \emph{vosotros}, comme s'il s'adressait à toute la famille : l'accroche, qu'on ne traduit pas, devait justement permettre de trancher. Autres fautes relevées : \emph{dormía} mal conjugué (verbe régulier à l'imparfait), \emph{como si} suivi d'un autre temps que le subjonctif imparfait, et des périphrases maladroites pour \emph{vomir} (\emph{vomitar}).

\end{piege}

\begin{jury}

Rapport 2026 (p.~57-58) : dans l'extrait, le narrateur s'adresse à son père, « c'est donc « usted » qu'il fallait utiliser ». Le jury attendait aussi \emph{los hay que \ldots{} otros que} (ou \emph{unos \ldots{} otros}), \emph{como si nos llamara} et \emph{yo dormía}, et rappelle que tout le texte doit être traduit.

\end{jury}

\hypertarget{partie-v-italien}{%
\section{Partie V --- Italien}\label{partie-v-italien}}

Article de presse sur une canicule précoce et ses causes marines. La question 2 comporte deux volets, \emph{sarà possibile ?} et \emph{in che modo ?}. Le thème est un monologue sur la dépendance au téléphone, plein de faux amis.

\hypertarget{question-13-italien-expression-uxe9crite-1.}{%
\subsection{Question 13 --- Italien, expression écrite 1.}\label{question-13-italien-expression-uxe9crite-1.}}

\begin{indication}{1}

\emph{Illustri} demande un exposé, pas un avis. Retenez trois aspects : l'anticipation sur la saison, l'ampleur dans les villes (jour et nuit) et la cause de fond.

\end{indication}

\begin{indication}{2}

Gardez deux ou trois chiffres frappants au plus, et nommez la cause que l'article présente comme le véritable moteur, avec sa conséquence : une mer qui accumule de l'énergie.

\end{indication}

\hypertarget{solution-12}{%
\subsubsection{Solution}\label{solution-12}}

Già a fine giugno, con largo anticipo sull'estate, l'anticiclone africano porta un caldo eccezionale nelle città: il ministero della Salute ha messo in bollino rosso diciassette città su ventisette, presto diciannove, e a Taranto si attendono 40 gradi. Neppure la notte dà sollievo, perché le minime restano sopra i 20 gradi e l'afa accresce il disagio. La causa profonda è il riscaldamento del Mediterraneo, molto più caldo della media: accumulando energia, il mare non mitiga più il caldo e rende possibili fenomeni violenti.

\textbf{(83 parole)}

\begin{piege}

\emph{il caldo} et \emph{la calura} désignent la chaleur subie, \emph{il calore} la chaleur au sens physique. Devant les degrés, l'italien met l'article : \emph{i 38 gradi}. \emph{Bollino rosso}, désignation propre au dispositif italien, se garde tel quel.

\end{piege}

\begin{jury}

Rapport 2026 (p.~60) : le jury dénonce la réponse sous forme de catalogue, une « tendance détestable » heureusement rare. La réponse doit synthétiser les idées du texte de façon organisée : « Le jury ne saurait se contenter d'un inventaire, fût-il complet. » D'où une réponse en trois temps (anticipation, villes jour et nuit, cause marine), et non une suite de chiffres.

\end{jury}

\hypertarget{question-14-italien-expression-uxe9crite-2.}{%
\subsection{Question 14 --- Italien, expression écrite 2.}\label{question-14-italien-expression-uxe9crite-2.}}

\begin{indication}{1}

Traitez les deux volets. Pour le premier, interrogez le verbe \emph{risolvere} : une crise climatique se « résout »-elle comme une équation ?

\end{indication}

\begin{indication}{2}

Pour le second, distinguez des leviers de nature différente (technique, politique, adaptation) et illustrez chacun par un fait vérifiable. Une introduction peut partir de votre propre expérience de la chaleur.

\end{indication}

\begin{indication}{3}

Pour exprimer votre opinion, \emph{Secondo me} ou \emph{A mio parere} + indicatif sont sûrs. Si vous employez \emph{credo che} / \emph{penso che} avec un présent, mettez le subjonctif.

\end{indication}

\hypertarget{solution-13}{%
\subsubsection{Solution}\label{solution-13}}

Le ondate di calore come quella di giugno 2025 lo dimostrano: chi, come molti studenti, vive in un alloggio senza aria condizionata sa che la crisi climatica è già qui. Si potrà risolverla un giorno? Secondo me, non come si risolve un'equazione, perché l'anidride carbonica resta per secoli nell'atmosfera. È però possibile stabilizzare il clima, a tre condizioni.

La prima è tecnica. In Francia quasi tutta l'elettricità proviene dal nucleare e dalle rinnovabili, mentre in Danimarca l'eolico supera ormai la metà della produzione. Per cemento, acciaio e aviazione serviranno idrogeno, cattura del carbonio e, semplicemente, meno consumi.

La seconda è politica, ed è la più dura. Un accordo internazionale funziona solo se nessuno ha interesse a tradirlo: per questo servono un prezzo del carbonio e un meccanismo di adeguamento alle frontiere, come quello adottato dall'Unione europea.

La terza è l'adattamento delle città descritte nell'articolo: piantare alberi, ridurre l'asfalto, rivedere gli orari di lavoro durante le ondate di calore.

Sono quindi moderatamente ottimista: non credo né alla catastrofe inevitabile né al miracolo tecnologico, ma a mille decisioni concrete, da prendere subito.

\textbf{(180 parole)}

\begin{methode}

Pour une question double, consacrez un paragraphe à chaque volet, ou annoncez dès la thèse comment vous allez répondre au second (\emph{a tre condizioni}) : le correcteur voit tout de suite que la consigne est respectée.

\end{methode}

\begin{piege}

Distinguez \emph{credere a} (croire à quelque chose, y ajouter foi) et \emph{credere in} (croire en une personne ou en une valeur). Veillez aussi au référent des pronoms : un \emph{stabilizzarla} renverrait à la crise, non au climat.

\end{piege}

\begin{jury}

Rapport 2026 (p.~60) : pour cette question, « on a pu regretter parfois un manque d'ouverture » ; une réflexion personnelle plus large, étayée d'exemples variés, était bienvenue. Le jury note que les candidats, souvent logés sans climatisation, peuvent nourrir leur réflexion de leur propre témoignage.

\end{jury}

\hypertarget{question-15-italien-thuxe8me}{%
\subsection{Question 15 --- Italien, thème}\label{question-15-italien-thuxe8me}}

\begin{indication}{1}

Repérez les faux amis : \emph{portable}, \emph{addiction}, \emph{agenda}. Demandez-vous aussi quel verbe convient pour un téléphone qui sonne.

\end{indication}

\begin{indication}{2}

La dernière phrase du troisième paragraphe enchaîne un conditionnel, puis deux subjonctifs. Déterminez le temps de chacun.

\end{indication}

\begin{indication}{3}

Pour \emph{Constance, c'est différent}, ne calquez pas la dislocation : on veut dire que la situation est différente, pas Constance elle-même.

\end{indication}

\hypertarget{solution-14}{%
\subsubsection{Solution}\label{solution-14}}

Abdil ha appena riattaccato che Constance mi chiama. Nelle mani di uno scorticato nell'anima, il telefono moltiplica gli SOS. Finché c'è qualcuno all'altro capo del filo, si conserva un po' di speranza. È sorprendente quanto io sia solo e quanto spesso mi chiamino {[}\ldots{]}.

Il cellulare è diventato un personaggio centrale della mia esistenza. È la mia dipendenza numero uno {[}\ldots{]}. Quando lo perdo, sono ancora più infelice del solito {[}\ldots{]}. A volte mi chiedo se non abiti nel mio cellulare. Ammiro quelli che ne fanno a meno. Vorrei essere come loro {[}\ldots{]}.

Quando squilla, vuol dire che esisto ancora un po'. Chiamo anch'io. Mi capita di consultare la mia lista dei contatti o una vecchia agenda per rifarmi vivo con un amico perso di vista. Possono essere passati vent'anni, ma volano così in fretta, vent'anni. Non vorrei che si spegnessero, i miei amici perduti, senza che ci fossimo detti prima qualche parola d'affetto {[}\ldots{]}.

Con Constance è diverso. Quando chiede aiuto cinque o sei volte al giorno, è una sofferenza.

Choix justifiés :

\begin{itemize}
\tightlist
\item
  \emph{À peine \ldots{} que} devient \emph{ha appena \ldots{} che}.
\item
  \emph{le portable} se dit \emph{il cellulare} ; \emph{il portatile} désigne l'ordinateur.
\item
  \emph{addiction} se dit \emph{dipendenza} ; \emph{s'en passer}, \emph{farne a meno}.
\item
  \emph{Je ne voudrais pas qu'ils s'éteignent \ldots{} sans qu'on se soit dit} se traduit par \emph{Non vorrei che si spegnessero \ldots{} senza che ci fossimo detti} : conditionnel présent, subjonctif imparfait, subjonctif plus-que-parfait.
\item
  \emph{È sorprendente quanto io sia solo} : le subjonctif suit l'exclamative indirecte, et le pronom \emph{io} lève l'ambiguïté de \emph{sia}.
\item
  \emph{squillare} est le verbe propre au téléphone, mais \emph{suonare} n'est pas fautif.
\end{itemize}

\begin{methode}

En italien, la concordance des temps suit le verbe principal. Après un conditionnel présent de volonté (\emph{vorrei}), la subordonnée se met au subjonctif imparfait, et l'antériorité au subjonctif plus-que-parfait.

\end{methode}

\begin{piege}

Fautes relevées par le jury sur des mots simples : \emph{malheureux} rendu par \emph{sfelice} ou \emph{maloroso} (on dit \emph{infelice}), \emph{étonnant} par \emph{surprenante} ou \emph{stuppifacente} (on dit \emph{sorprendente}). Sur les verbes irréguliers, \emph{volerebbe} confond \emph{volere} et \emph{volare} ; le subjonctif imparfait de \emph{volere} est \emph{volessi}, \emph{volesse}. \emph{Il m'arrive de} se rend par \emph{capitare}.

\end{piege}

\begin{jury}

Rapport 2026 (p.~59-60) : « Une petite difficulté a été l'emploi du verbe à l'imparfait du Subjonctif en italien après un conditionnel dans la principale », dans \emph{Je ne voudrais pas qu'ils s'éteignent}. Le jury relève aussi \emph{venir de} + infinitif, des omissions sanctionnées et des temps non respectés (présent au lieu du conditionnel ou du subjonctif imparfait), et recommande une relecture attentive du texte et de la traduction en regard.

\end{jury}

\hypertarget{partie-vi-russe}{%
\section{Partie VI --- Russe}\label{partie-vi-russe}}

L'article décrit les codes vestimentaires et l'argot d'une culture jeune venue du hip-hop. La question 2 appelle une réponse personnelle et nuancée. Le thème, au présent de narration avec des futurs d'anticipation, oblige à choisir l'aspect verbe par verbe.

\hypertarget{question-16-russe-expression-uxe9crite-1.}{%
\subsection{Question 16 --- Russe, expression écrite 1.}\label{question-16-russe-expression-uxe9crite-1.}}

\begin{indication}{1}

Le texte décrit un style, et non spécifiquement « la jeunesse russe ». Répondez à partir du texte en organisant l'inventaire en trois blocs : vêtements, marques corporelles, langage.

\end{indication}

\begin{indication}{2}

La dernière phrase doit dire la fonction de ce style selon l'auteur : reconnaissance, protestation, mise à distance des adultes.

\end{indication}

\hypertarget{solution-15}{%
\subsubsection{Solution}\label{solution-15}}

Судя по тексту, современная молодёжь выделяется прежде всего свободным, как правило спортивным стилем, пришедшим из американского хип-хопа и рэпа. Молодые люди носят рваные джинсы, незашнурованные кроссовки, очень широкие брюки, худи и футболки на несколько размеров больше, яркие золотистые украшения на шее, солнечные очки и бейсболки козырьком набок. К этому добавляются татуировки и особый сленг: «лол», «рофл», «шеймить». Все эти знаки позволяют молодым людям узнавать «своих», выражать протест и отделять себя от мира взрослых. Это одновременно мода, вызов и способ самовыражения.

\textbf{(80 слов)}

\begin{piege}

\emph{молодёжь} est un collectif singulier féminin : \emph{молодёжь отличается}, et non \emph{отличаются}. Le génitif pluriel suit les quantificateurs : \emph{на несколько размеров}.

\end{piege}

\begin{jury}

Rapport 2026 (p.~61) : dans cette question, « certains candidats ont oublié des éléments importants ». D'où l'inventaire en trois blocs, qui n'oublie ni l'origine hip-hop ni la fonction du style. Le jury relève aussi une orthographe phonétique et une ponctuation approximative, alors que « les règles de ponctuation du russe sont strictes ».

\end{jury}

\hypertarget{question-17-russe-expression-uxe9crite-2.}{%
\subsection{Question 17 --- Russe, expression écrite 2.}\label{question-17-russe-expression-uxe9crite-2.}}

\begin{indication}{1}

\emph{В какой мере} demande une réponse graduée : distinguez les domaines où vous vous reconnaissez de ceux où vous ne vous reconnaissez pas.

\end{indication}

\begin{indication}{2}

Partez de détails concrets de votre vie d'étudiant, puis remontez à l'idée générale du texte (\emph{мой язык --- только для своих}) pour la discuter.

\end{indication}

\begin{indication}{3}

Pour un lecteur russe, précisez ce qu'est une classe préparatoire : \emph{подготовительный класс} seul évoque l'année avant l'école primaire.

\end{indication}

\hypertarget{solution-16}{%
\subsubsection{Solution}\label{solution-16}}

Я узнаю себя в этом стиле лишь отчасти.

Внешне --- почти нет. Я учусь в подготовительном классе, где готовятся к вступительным экзаменам в инженерные школы, и мой «стиль» определяется расписанием: удобные брюки, свитер, рюкзак с книгами. У меня нет ни татуировок, ни золотых цепочек, а бейсболку я ношу козырьком вперёд. Правда, кроссовки у меня иногда остаются незашнурованными, но только из-за спешки, а не из протеста. Мне просто некогда думать о моде: после десяти часов занятий важнее выспаться.

Зато в языке я узнаю себя сразу. Мы с друзьями тоже говорим «лол» и «рофл», «шеймим» и «прикалываемся», и родители часто не понимают половины наших разговоров. Кроме того, у нас есть свой школьный жаргон, который никто со стороны не поймёт. Это тот самый механизм, о котором пишет автор: язык только для своих.

Поэтому мне кажется, что автор прав в главном. Молодёжная культура --- это не набор вещей, а способ провести границу между «своими» и «чужими». У рэпера это рваные джинсы, у студента --- сокращения, у геймера --- никнейм. Каждое поколение придумывает свои знаки: у моих родителей были другие слова и другая музыка. Форма меняется, а потребность сказать «я --- свой» остаётся.

\textbf{(183 слова)}

\begin{methode}

Le russe exige la double négation (\emph{никто \ldots{} не поймёт}, \emph{ни татуировок, ни цепочек}). Après \emph{нет}, on emploie le génitif ; après \emph{оставаться}, l'attribut se met à l'instrumental (\emph{остаются незашнурованными}).

\end{methode}

\begin{piege}

Le nombre de mots s'accorde lui aussi : le nom suit le dernier chiffre. \emph{1 слово}, \emph{2, 3, 4 слова}, \emph{5 à 20 слов} : on écrit donc \emph{80 слов}, mais \emph{183 слова} et \emph{181 слово}.

\end{piege}

\begin{jury}

Rapport 2026 (p.~61) : l'essai a parfois donné des compositions ne répondant pas à la question, « hors sujet, donc », ou écrites au fil de la plume sans plan construit. Le nombre de mots n'était pas toujours respecté. Les notes vont de 10 à 16, pour une moyenne de 13,25.

\end{jury}

\hypertarget{question-18-russe-thuxe8me}{%
\subsection{Question 18 --- Russe, thème}\label{question-18-russe-thuxe8me}}

\begin{indication}{1}

Reconstituez la généalogie avant de traduire : qui est la « belle-mère » de la mère du père ? Le russe a des mots différents selon le côté de la famille.

\end{indication}

\begin{indication}{2}

Le présent de narration se rend par l'imperfectif présent. Les futurs (\emph{finira par}, \emph{ne quittera jamais}) et le futur antérieur de bilan (\emph{il aura fallu}) demandent une décision d'aspect.

\end{indication}

\begin{indication}{3}

\emph{le père haï autant qu'espéré} est très ramassé : cherchez une construction comparative russe qui garde ce parallélisme sans changer de sujet.

\end{indication}

\hypertarget{solution-17}{%
\subsubsection{Solution}\label{solution-17}}

Мой отец очень близок со своей бабушкой. Она рядом с ним на протяжении всего его детства и отрочества, проведённых вместе с братьями в маленьком доме в Ле-Везине. Их мать работает в Париже и в конце концов там поселится --- быть может, не выдержав властного характера свекрови. Анета никогда не расстанется со своим сыном Дато; она ему предана, и он отвечает ей тем же, защищая её от Иосифа --- отца, столь же ненавистного, сколь и долгожданного. Она умирает, когда моя мать беременна мной, с нетерпением ожидая ребёнка, которого ей не суждено узнать и которого она уже называет «девочкой». Иосиф в своих записях ни разу не упоминает ни Анету, ни какую-либо другую женщину или ребёнка, словно прожил жизнь в полном одиночестве.

Понадобилось целое стечение обстоятельств, много лет спустя после смерти отца, чтобы я задумалась о своём прадеде по отцовской линии и об этой стороне своего происхождения.

Choix justifiés :

\begin{itemize}
\tightlist
\item
  La belle-mère de la mère du père est la mère de son mari : \emph{свекровь} (\emph{тёща} et \emph{мачеха} fausseraient la généalogie).
\item
  Les futurs d'anticipation sont rendus par des perfectifs futurs (\emph{поселится}, \emph{не расстанется}). Le futur antérieur de bilan devient \emph{Понадобилось \ldots, чтобы я задумалась} : \emph{чтобы} + passé, car les sujets diffèrent.
\item
  \emph{haï autant qu'espéré} devient \emph{столь же ненавистного, сколь и долгожданного}.
\item
  Le gérondif \emph{защищая} a pour sujet celui de la proposition principale (Dato), et \emph{ожидая} celui de sa propre principale (Anéta).
\item
  On transcrit les noms (\emph{Анета}, \emph{Дато}, \emph{Иосиф}, \emph{Ле-Везине}). Les références et le titre, eux, ne se traduisent pas.
\end{itemize}

\begin{piege}

\emph{беременна мной} est la tournure naturelle pour « enceinte de moi ». \emph{ждёт меня} est plus ambigu, surtout à côté de \emph{ожидая ребёнка}.

\end{piege}

\begin{jury}

Rapport 2026 (p.~61) : le vocabulaire du thème était quotidien (famille, maison), et le russe des candidats souvent riche ; c'est la compréhension du français qui a parfois fait défaut, d'où des erreurs et des « oublis ». Reconstituer la généalogie et le sens de \emph{haï autant qu'espéré} avant d'écrire évite l'essentiel de ces pertes.

\end{jury}

\hypertarget{bilan}{%
\section{Bilan}\label{bilan}}

\textbf{À retenir pour toutes les langues.}

\begin{itemize}
\tightlist
\item
  Pour la question 1 : répondre directement, sans introduction ; relever quatre ou cinq éléments, les regrouper avec des connecteurs, reformuler, sans rien ajouter ni retrancher au sens du texte.
\item
  Pour la question 2 : une ou deux phrases d'introduction qui situent le sujet dans l'actualité et posent le problème, une prise de position nette et nuancée, deux ou trois paragraphes chacun avec son idée et son exemple, une conclusion qui répond. Les exemples viennent de vos connaissances, pas de l'article.
\end{itemize}

\textbf{Ce que le jury 2026 a jugé discriminant.} En anglais, les deux « depuis » du thème et la richesse de langue de l'essai ; en espagnol, le \emph{usted} du thème ; en arabe, la structure en paragraphes et le nombre de mots ; partout, les exemples personnels et la reformulation en question 1.
- Au thème : fixer d'abord le repère temporel, traduire la phrase et non le mot, préférer l'idiome au calque, et ne rien laisser en blanc.
- Compter ses mots et en indiquer le nombre. Viser le milieu de la fourchette laisse une marge si les conventions de décompte diffèrent (élisions, nombres, traits d'union).

\textbf{Grammaire à revoir, langue par langue.} Allemand : place du verbe et \emph{Konjunktiv I}. Anglais : \emph{past perfect} avec \emph{for}, \emph{present perfect} avec \emph{since}, \emph{want someone to do} et passif après \emph{want}, prétérit des verbes irréguliers. Arabe : duel, accord du verbe placé avant son sujet, négations (لم، لن، لا). Espagnol : \emph{usted} et \emph{vosotros}, subjonctif imparfait après \emph{como si}, \emph{sino} après une négation, opposition \emph{ser} / \emph{estar}. Italien : concordance des temps au subjonctif, faux amis. Russe : aspect, génitif après une négation, instrumental attribut.

\textbf{Situations qui se retrouvent dans d'autres sujets.} Expliquer les difficultés d'une entreprise ou d'un groupe à partir d'un article chiffré. Discuter un critère unique (le prix, la technique, l'idéologie). Exposer des politiques publiques avec des exemples nationaux. Traduire un dialogue familier inséré dans un récit au passé.

\end{document}
